% block: matrix-function-definition-heading page:203 type:section
\subsection{{矩阵函数的定义与性质}}\label{sec:matrix-function-properties}

% block: matrix-spectrum-introduction page:203 type:ordinary-paragraph
{矩阵 $A$ 的全部特征值构成的集合称为 $A$ 的谱. 矩阵的谱上的函数定义如下.}

% block: function-on-spectrum-definition page:203 type:definition
\Needspace{8\baselineskip}
\begin{Definition}{{矩阵的谱上的函数}}\label{def:function-on-matrix-spectrum}
{设 $\lambda_1,\lambda_2,\ldots,\lambda_s$ 是矩阵 $A$ 的互异特征值, $n_i$ 是对应于 $A$ 的特征值 $\lambda_i$ 的一个 Jordan 块的大小. 如果存在下列值:}
\[f(\lambda_i),\frac{\mathrm{d}}{\mathrm{d}x}f(\lambda_i),\ldots,\frac{\mathrm{d}^{n_i-1}}{\mathrm{d}x^{n_i-1}}f(\lambda_i)\quad(i=1,2,\ldots,s),\]
{则称 $f(x)$ 为矩阵 $A$ 的谱上的函数.}
\end{Definition}

% block: matrix-function-derivative-notation page:203 type:ordinary-paragraph
{以下用符号 $f^{(i)}(x)$ 代替 $\frac{\mathrm{d}^i}{\mathrm{d}x^i}f(x)$.}

% block: inverse-function-spectrum-example page:203 type:example
\Needspace{7\baselineskip}
\begin{Example}\label{ex:inverse-on-jordan-spectrum}
{设 $A$ 为 Example~\ref{ex:jordan-block-diagonal-six} 中的矩阵 $J$. 由于 $\lambda_1=\lambda_2=1$ 且 $\lambda_3=2$, 互异特征值的个数为 $s=2$. 对应于 $\lambda_1=\lambda_2=1$ 和 $\lambda_3=2$ 的 Jordan 块的最大大小分别为 $n_1=2$ 和 $n_2=3$. 矩阵 $A$ 的谱为 $\{\lambda_1,\lambda_3\}$. 现在给出矩阵 $A$ 的谱上的函数的一个例子. 设 $f(x)=x^{-1}$. 则 $f(x)$ 是矩阵 $A$ 的谱上的函数, 因为 $f^{(1)}(x)=-x^{-2}$, $f^{(2)}(x)=2x^{-3}$, 且存在下列谱上的函数值: $f(\lambda_1)=1$, $f^{(1)}(\lambda_1)=-1$, $f(\lambda_3)=1/2$, $f^{(1)}(\lambda_3)=-1/4$ 及 $f^{(2)}(\lambda_3)=1/4$.}
\end{Example}

% block: jordan-matrix-function-transition page:203 type:ordinary-paragraph
{利用 Theorem~\ref{thm:jordan-canonical-form} 中的 Jordan 标准形和 Definition~\ref{def:function-on-matrix-spectrum} 中矩阵的谱上的函数, 可以如下定义 $A$ 的矩阵函数:}

% block: jordan-matrix-function-definition page:204 type:definition
\Needspace{9\baselineskip}
\begin{Definition}{{矩阵函数}}\label{def:jordan-matrix-function}
{设 $f(x)$ 是矩阵 $A\in\mathbb{C}^{n\times n}$ 的谱上的函数, $A=ZJZ^{-1}$ 是 $A$ 的 Jordan 标准形. 则矩阵函数 $f(A)=f(ZJZ^{-1})$ 定义如下:}
\[f(A)=Z\operatorname{diag}(f(J_1),f(J_2),\ldots,f(J_p))Z^{-1}.\]
{这里 $f(J_k)$ 为下面的 $m_k\times m_k$ 矩阵:}
\[
f(J_k)=\begin{bmatrix}f(\lambda_k)&\frac{f^{(1)}(\lambda_k)}{1!}&\cdots&\frac{f^{(m_k-1)}(\lambda_k)}{(m_k-1)!}\\&f(\lambda_k)&\ddots&\vdots\\&&\ddots&\frac{f^{(1)}(\lambda_k)}{1!}\\&&&f(\lambda_k)\end{bmatrix}\in\mathbb{C}^{m_k\times m_k},
\]
{且 $m_k\times m_k$ 是对应于特征值 $\lambda_k$ 的 Jordan 块的大小.}
\end{Definition}

% block: matrix-function-eigenvalues page:204 type:ordinary-paragraph
{由 Definition~\ref{def:jordan-matrix-function} 可知, 当 $\lambda_k$ 是矩阵 $A$ 的一个特征值时, 矩阵函数 $f(A)$ 的特征值为 $f(\lambda_k)$.}

% block: matrix-function-diagonalizable-case page:204 type:ordinary-paragraph
{特别地, 当矩阵 $A$ 可对角化时, Jordan 块为对角矩阵 $D_k$, $k=1,2,\ldots,p$. 于是由 $f(A)=Z\operatorname{diag}(D_1,D_2,\ldots,D_p)Z^{-1}$ 可知, $f(A)$ 与 $A$ 具有相同的特征向量.}

% block: inverse-matrix-function-example page:204 type:example
\Needspace{9\baselineskip}
\begin{Example}\label{ex:inverse-jordan-matrix-function}
{设 $A$ 为 Example~\ref{ex:jordan-block-diagonal-six} 中的矩阵 $J$, 并设 $f(x)=x^{-1}$. $A$ 的 Jordan 标准形为 $A=ZJZ^{-1}$, 其中 $Z$ 是单位矩阵. 则 $f(A)=f(ZJZ^{-1})=f(J)$, 因而由 $\lambda_1=\lambda_2=1$, $\lambda_3=2$ 及 Definition~\ref{def:jordan-matrix-function}, 有}
\[
\begin{aligned}
f(A)&=\left[\begin{array}{cc|c|ccc}f(\lambda_1)&\frac{f^{(1)}(\lambda_1)}{1!}&0&0&0&0\\0&f(\lambda_1)&0&0&0&0\\\hline0&0&f(\lambda_2)&0&0&0\\\hline0&0&0&f(\lambda_3)&\frac{f^{(1)}(\lambda_3)}{1!}&\frac{f^{(2)}(\lambda_3)}{2!}\\0&0&0&0&f(\lambda_3)&\frac{f^{(1)}(\lambda_3)}{1!}\\0&0&0&0&0&f(\lambda_3)\end{array}\right]\\
&=\left[\begin{array}{cc|c|ccc}1&-1&0&0&0&0\\0&1&0&0&0&0\\\hline0&0&1&0&0&0\\\hline0&0&0&2^{-1}&-4^{-1}&8^{-1}\\0&0&0&0&2^{-1}&-4^{-1}\\0&0&0&0&0&2^{-1}\end{array}\right].
\end{aligned}
\]
{容易看出 $f(A)$ 是矩阵 $A$ 的逆, 即从矩阵函数 $f(A)$ 的定义导出了 $A^{-1}$.}
\end{Example}

% block: square-root-exponential-jordan-example page:205 type:example
\Needspace{8\baselineskip}
\begin{Example}\label{ex:square-root-exponential-jordan}
{我们考虑 $A$ 的矩阵函数的其他例子. 设}
\[A=\begin{bmatrix}-1&6&-9\\4&-4&12\\3&-5&11\end{bmatrix}.\]
{则 Jordan 标准形为}
\[A=ZJZ^{-1},\]
{其中}
\[Z^{-1}AZ=J=\begin{bmatrix}2&1&0\\0&2&1\\0&0&2\end{bmatrix}\]
{且}
\[Z=\begin{bmatrix}3&0&-1\\0&2&1\\-1&1&1\end{bmatrix},\quad Z^{-1}=\begin{bmatrix}1&-1&2\\-1&2&-3\\2&-3&6\end{bmatrix}.\]
{利用这些矩阵, 对应于 $f(x)=x^{1/2}$ 的 $A^{1/2}$ 为}
\[A^{1/2}=Z\begin{bmatrix}\sqrt{2}&\frac{1}{2\sqrt{2}}&-\frac{1}{16\sqrt{2}}\\0&\sqrt{2}&\frac{1}{2\sqrt{2}}\\0&0&\sqrt{2}\end{bmatrix}Z^{-1}.\]
{容易看出 $A^{1/2}A^{1/2}=A$.}

{最后一个例子是 $A$ 的矩阵指数函数. 对应于 $f(x)=e^x$ 的矩阵指数函数 $e^A$ 为}
\[e^A=Z\begin{bmatrix}e^2&e^2&\frac{e^2}{2}\\0&e^2&e^2\\0&0&e^2\end{bmatrix}Z^{-1}.\]
{注意, Jordan 标准形的计算在数值上是不稳定的. 因此, 这些例子并不用于实际计算.}
\end{Example}

% block: matrix-function-properties-transition page:205 type:ordinary-paragraph
{下面列举矩阵函数的一些性质. 详情可参见例如 \BookCite{97}.}

% block: matrix-function-properties-theorem pages:205,206 type:theorem
\Needspace{14\baselineskip}
\begin{Theorem}{{矩阵函数的性质}}\label{thm:matrix-function-properties}
{设 $f(x)$ 是矩阵 $A$ 的谱上的函数. 则:}
\begin{enumerate}[label=(\arabic*)]
\item $Af(A)=f(A)A$.
\item $XA=AX\Rightarrow Xf(A)=f(A)X$.
\item $f(A^{\top})=f(A)^{\top}$.
\item $f(XAX^{-1})=Xf(A)X^{-1}$.
\item {如果 $\lambda_k$ 是 $A$ 的一个特征值, 则 $f(\lambda_k)$ 是 $f(A)$ 的一个特征值.}
\end{enumerate}
{设 $f(x)$ 与 $g(x)$ 是矩阵 $A$ 的谱上的函数. 则:}
\begin{enumerate}[label=(\arabic*),start=6]
\item {如果 $(f+g)(x):=f(x)+g(x)$, 则 $(f+g)(A)=f(A)+g(A)$.}
\item {如果 $(fg)(x):=f(x)g(x)$, 则 $(fg)(A)=f(A)g(A)$.}
\end{enumerate}
{利用 Definition~\ref{def:function-on-matrix-spectrum}, 可得}
\begin{enumerate}[label=(\arabic*),start=8]
\item {$f(A)=g(A)$ 当且仅当\newline{}$f(\lambda_i)=g(\lambda_i),\ldots,f^{(n_i-1)}(\lambda_i)=g^{(n_i-1)}(\lambda_i)$\newline{}$(i=1,2,\ldots,s)$.}
\end{enumerate}
\end{Theorem}

% block: jordan-matrix-function-computational-warning page:206 type:ordinary-paragraph
{\leavevmode\phantomsection\label{srcpage:printed-194}由 Definition~\ref{def:jordan-matrix-function}, 可以定义与初等函数 $x^{1/n}$, $e^x$, $\sin x$, $\cos x$, $\log x$ 对应的 $A^{1/n}$, $e^A$, $\sin A$, $\cos A$, $\log A$. 然而, 由于 Jordan 标准形的数值计算相当不稳定, 不建议通过 Jordan 标准形计算矩阵函数. 因此, 许多研究者设计了数值稳定的矩阵函数算法.}

% block: matrix-functions-numerical-algorithms-transition page:206 type:ordinary-paragraph
{以下介绍计算矩阵函数的一些数值算法. 我们将看到, Krylov 子空间的概念有助于计算其中一些矩阵函数.}
